\documentclass[11pt]{article}
\usepackage[a4paper, margin=2.0cm]{geometry}
\usepackage[T1]{fontenc}
\usepackage[utf8]{inputenc}
\usepackage{booktabs, array, amsmath, amssymb, xcolor, enumitem}
\usepackage[colorlinks=true, linkcolor=blue!50!black, urlcolor=blue!50!black]{hyperref}
\newcolumntype{L}[1]{>{\raggedright\arraybackslash}p{#1}}
\setlength{\parskip}{4pt}
\sloppy

% Every number below comes from experiments/openfhe/standard_tables/ (generated
% by experiments/openfhe/make_standard_tables.py from aggregate.py's outputs).
\newcommand{\tab}[1]{../../experiments/openfhe/standard_tables/#1}
\input{\tab{numbers.tex}}
\newcommand{\R}[1]{\csname r:#1\endcsname}   % ratio OpenFHE/SEAL, e.g. \R{CKKS:8192:1:relinearize}
\newcommand{\X}[1]{\csname x:#1\endcsname}   % its inverse (SEAL/OpenFHE)
\newcommand{\V}[1]{\csname v:#1\endcsname}   % mean latency, ms, e.g. \V{OpenFHE:CKKS:16384:1:add}
\newcommand{\RM}[1]{\csname rm:#1\endcsname} % ratio of medians
\newcommand{\RQ}[1]{\csname rq:#1\endcsname} % quartile envelope of the ratio

\title{OpenFHE Standard campaign: results\\[2pt]
\large SEAL 4.1.2 and OpenFHE 1.5.1, same sessions, \nConfigs\ configurations\\ (version 2)}
\author{Boungo Tr\'esor}
\date{\today}

\begin{document}
\maketitle

\noindent\textbf{Scope.} This note reports the complete Standard campaign authorized after the pilot (\texttt{openfhe\_pilot\_note}): every eligible configuration, OpenFHE and SEAL measured in the same sessions with the pilot's harness and protocol. It gives the consolidated results, the excluded and problematic cells, and an interpretation of the main differences. Latency only: energy and memory are a later step (Section~\ref{sec:limits}). All tables and numbers are generated from the data (\texttt{experiments/openfhe/make\_standard\_tables.py}); every incident is recorded in \texttt{experiments/results/raw/standard\_campaign/INCIDENTS.md}.

\paragraph{Changes from version 1} (tag \texttt{openfhe-standard-v1}, unchanged) in answer to its review: every conclusion is stated with the configurations and boundaries that support it, and no general statement about one library being faster is made (Section~\ref{sec:interp}); key generation is described as an initialization and key-generation cost; the sensitivity of each ratio to the estimator (means, medians, quartiles) and the reproducibility of the main ratios across independent sessions are reported, and the meaning of the confidence intervals is stated (Section~\ref{sec:stats}). The CKKS 16384/1 relinearization pair, affected by a kernel change, will be measured again under one kernel (Section~\ref{sec:problems}); its current values are kept here and marked.

\paragraph{Summary.}
\begin{itemize}[leftmargin=*, itemsep=2pt]
  \item \textbf{\nCells\ valid cells in \nSessions\ sessions enter the analysis (\nExcluded\ excluded).} Every one passed the fail-closed validation: complete repetitions, every OpenFHE result checked correct, configuration card identical to the session manifest.
  \item \textbf{\nSetAside\ earlier measurements are not part of these \nCells.} They had been measured while the laptop ran on battery, which lowers the CPU frequency; they are kept in the raw data, set aside, and their cells were measured again on mains. A power guard then prevented any measurement on battery.
  \item \textbf{Main observations} (matched-profile comparisons of equivalent boundaries, same session; ratio = OpenFHE/SEAL). Initialization and key generation: ratio \ratioMinkeygen--\ratioMaxkeygen\ at CKKS 2048/5, 8192/1 and 16384/1. Decryption to SEAL's boundary: \ratioMinckksdecrypt--\ratioMaxckksdecrypt\ in the \ratioNckksdecrypt\ CKKS configurations. Multiplication without relinearization: \ratioMinmultiply--\ratioMaxmultiply\ at CKKS 4096/1, 8192/1, 16384/1. Key switching at CKKS 8192/1 and 16384/1: relinearization \ratioMinrelinearize--\ratioMaxrelinearize, rotation \ratioMinrotate\ (8192/1). At CKKS 8192/1, the only configuration measured in several independent sessions, these ratios were reproduced in \nRepSessions\ sessions (Table~\ref{tab:repro}).
\end{itemize}

\section{Protocol and sessions}

The protocol is the pilot's, unchanged (pilot note, Section~3): one process per cell; fresh inputs for every call; 100 measured repetitions after 5 warm-up, inner loop 1000 (keygen and rotate: one call); one pinned core, governor \texttt{performance}; both libraries in the same session in randomized order (seed recorded); laptop otherwise idle. A session was resumed after an interruption only with an identical manifest (code commit, binary and image hashes, configuration files), so the completed cells were reused and the interrupted cell measured again from the start.

\begin{table}[h!]\centering\footnotesize\setlength{\tabcolsep}{4pt}
\input{\tab{sessions.tex}}
\caption{Sessions. Freq.: median effective core frequency over the session. SEAL/archived: SEAL's mean in this session over its mean in the archived SEAL campaign, same cell: minimum--median--maximum (number of cells). Data: \texttt{experiments/results/final/openfhe\_standard\_sessions.csv}, \texttt{openfhe\_standard.csv}.}
\label{tab:sessions}
\end{table}

SEAL ran faster in these sessions than in the archived SEAL campaign (median ratio \archMedian\ over \archN\ cells, range \archMin--\archMax). This is a difference between the two campaigns' conditions; it does not affect the OpenFHE/SEAL comparisons, which only use cells of the same session. The lowest ratio (\archMin) is SEAL's \emph{add} at BFV 16384/1, whose archived value is itself flagged HIGH\_VARIANCE.

\section{Results}

Results are reported in three groups, as in the pilot: (i) \emph{matched-profile comparisons} of \emph{equivalent computational boundaries}, the only cells where a ratio OpenFHE/SEAL is formed; (ii) equivalent boundaries of configurations that require a \emph{separate analysis} (the two libraries' moduli differ), reported per library without a ratio; (iii) \emph{end-to-end API costs}, each library's own call, which do not perform the same computation, without a ratio. No cell is an exact-parameter comparison. The class of every cell and its remaining differences are in \texttt{openfhe\_standard\_comparisons.csv}.

\begin{table}[h!]\centering\footnotesize\setlength{\tabcolsep}{4pt}
\input{\tab{matched.tex}}
\caption{Matched-profile comparisons, equivalent boundaries (\nRatios\ cells), same session. Ratio O/S of means, of medians, and quartile envelope (OpenFHE $Q_1$/SEAL $Q_3$ to OpenFHE $Q_3$/SEAL $Q_1$). Welch 90\% CI of the difference of means: within-cell variability only (Section~\ref{sec:stats}). decrypt$^\ast$: up to SEAL's decrypt boundary. Remaining differences (primes not equal, error distribution, key-switching method, CKKS scale where it differs) are listed per cell in \texttt{openfhe\_standard\_comparisons.csv}.}
\label{tab:matched}
\end{table}

\begin{table}[h!]\centering\scriptsize\setlength{\tabcolsep}{3pt}
\input{\tab{separate.tex}}
\caption{Equivalent boundaries of configurations requiring separate analysis (\nSeparateCells\ cells): mean latency per library, same session. The ciphertext moduli differ between the libraries, so no ratio is formed.}
\label{tab:separate}
\end{table}

\begin{table}[h!]\centering\footnotesize\setlength{\tabcolsep}{4pt}
\input{\tab{api.tex}}
\caption{End-to-end API costs (mean latency, ms), same session: each library's own call. \texttt{EvalAdd} copies its operand first (\texttt{EvalAddInPlace} does not); OpenFHE's \texttt{Decrypt} also decodes, SEAL's decrypt stops before; SEAL's encryption includes its special-prime step. The calls of a row do not perform the same computation and are not compared by a ratio.}
\label{tab:api}
\end{table}

\section{Statistical interpretation and reproducibility}
\label{sec:stats}

\paragraph{What the confidence intervals describe.} The 95\% CI of each mean and the Welch 90\% CI of each difference are computed from the 100 repetitions of one cell. The repetitions are successive, in one process, on one day; they are not independent experiments. These intervals therefore describe the variability \emph{within} a cell. They do not account for what changes \emph{between} sessions (process start, memory layout, frequency and thermal state of the day), and they are not used as a statement on the reproducibility of a result. That is assessed separately, below.

\paragraph{Sensitivity to the estimator.} For all but one of the \nRatios\ matched-profile ratios, the ratio of medians differs from the ratio of means by at most \medianDevMaxOthers\%, and the quartile envelope is narrow (Table~\ref{tab:matched}). The exception is the CKKS 8192/1 rotation: \R{CKKS:8192:1:rotate} from the means, \RM{CKKS:8192:1:rotate} from the medians, quartile envelope \RQ{CKKS:8192:1:rotate}. SEAL's rotation cell has a bimodal series: about 20 repetitions at its start and 18 at its end are slower (2.0--2.8~ms) than the stable middle (about 1.85~ms). It was the first cell of session s2: it ran for about one second, on a CPU that had just started, and the frequency record of that part of s2 was lost (logging bug, incident~4). Its OpenFHE partner ran in the middle of the session.

\paragraph{Reproducibility across independent sessions.} Only CKKS 8192/1 was measured, with both libraries, in more than one independent session: the two pilot runs (one with the laptop in use, kept as a record) and campaign session s2, on different days, each with new processes and its own randomized order. Table~\ref{tab:repro} gives the ratio in each. Between sessions, the relinearization ratio varies by less than 3\% (\repMinrelinearize--\repMaxrelinearize), decryption by about 10\% (\repMindecrypt--\repMaxdecrypt), multiplication by less than 8\% (\repMinmultiply--\repMaxmultiply) and initialization with key generation by about 20\% (\repMinkeygen--\repMaxkeygen); the highest values of the last three come from pilot run 1, made with the laptop in use. None of these variations changes the direction or the order of magnitude of a ratio. The rotation ratio varies most (\repMinrotate--\repMaxrotate\ from means, \repMedMinrotate--\repMedMaxrotate\ from medians), but OpenFHE's rotation is about twice SEAL's in every session.

\begin{table}[h!]\centering\footnotesize\setlength{\tabcolsep}{5pt}
\input{\tab{reproducibility.tex}}
\caption{CKKS 8192/1: ratio OpenFHE/SEAL of means (of medians in brackets) in each independent session that measured both libraries. Data: \texttt{experiments/results/final/openfhe\_standard\_reproducibility.csv}.}
\label{tab:repro}
\end{table}

\paragraph{Consequence for the rotation.} The existing observations support a stable conclusion at the level at which it is stated: at CKKS 8192/1, OpenFHE's rotation (one step, one Galois key) took about twice SEAL's time in all three sessions, with the exact factor uncertain by about $\pm$10\% (\repMinrotate--\repMaxrotate). A further measurement is therefore not needed for this conclusion; a more precise factor would need SEAL's cell to be measured again warm, in alternation with OpenFHE's.

For the other configurations, one session each was measured; the reproducibility of their ratios rests on the consistency across configurations (e.g.\ CKKS decryption \ratioMinckksdecrypt--\ratioMaxckksdecrypt\ in \ratioNckksdecrypt\ configurations measured in three different sessions) and on SEAL's agreement with the archived SEAL campaign (Table~\ref{tab:sessions}), not on repeated sessions of the same configuration.

\clearpage
\section{Excluded and problematic cells}
\label{sec:problems}

\paragraph{Excluded cells: none.} Validation excluded \nExcluded\ of \nCells\ cells (\texttt{openfhe\_standard\_excluded.csv} is empty).

\paragraph{Cells set aside: \nSetAside, all measured again (Table~\ref{tab:aside}).} On this laptop, running on battery makes the firmware lower the CPU power limit; the core frequency drops (to about 4.23--4.32~GHz against 4.40--4.46~GHz on mains in session s4). In s4, one cell (26~October) and eight cells (27~October, morning) ran partly or wholly on battery. They passed validation, but their timings come from a different power state, so they were moved, with their checks and cards, to \texttt{s4/set\_aside\_battery\_*} (not read by the aggregation) and measured again on mains with the same manifest. From then on every run went through a power guard (\texttt{experiments/tools/power\_guard.sh}): it records the mains state every 5~s, stops the measurement after 10~s on battery and resumes automatically 2~minutes after mains returns. Its record shows that every cell measured under it ran on mains.

\begin{table}[h!]\centering\footnotesize\setlength{\tabcolsep}{4pt}
\input{\tab{set_aside.tex}}
\caption{Cells measured on battery, set aside (kept in the raw data) and measured again on mains.}
\label{tab:aside}
\end{table}

\paragraph{High variance: \nHighVar\ of \nCells\ cells (Table~\ref{tab:hv}).} All are short operations (additions, CKKS 2048 decryption, rotations timed one call per repetition, full decryption) except OpenFHE's BFV 16384/1 key generation, whose flag comes from 3 isolated slow repetitions out of 100. The table gives the median and interquartile range next to the mean; except SEAL's CKKS 8192/1 rotation, where about a quarter of the repetitions are slower (upper quartile 2.18 against median 1.87~ms), the interquartile range is narrow, so the spread comes from a few outlying repetitions. All are reported as measured, with the flag; only one of them (SEAL's CKKS 8192/1 rotation) enters a ratio.

\begin{table}[h!]\centering\footnotesize\setlength{\tabcolsep}{4pt}
\input{\tab{highvar.tex}}
\caption{Cells above 5\% coefficient of variation (HIGH\_VARIANCE).}
\label{tab:hv}
\end{table}

\paragraph{Machine state between sessions.} With the same governor and power profile, the sustained frequency differed between days (s4: \freqMinsfour--\freqMaxsfour~GHz per cell; s5: \freqMinsfive--\freqMaxsfive~GHz; s1--s3: \freqMinsone--\freqMaxsone~GHz), for a reason not identified. Per-cell frequency, temperature and kernel are recorded in \texttt{openfhe\_standard.csv} and shown for s4 and s5 in Table~\ref{tab:cond}. The pairs used for a ratio in s5 ran within about 2\% of each other in frequency, except key generation (short cells, which end before the CPU heats up). Other points, all recorded in the incident log:
\begin{itemize}[leftmargin=*, itemsep=2pt]
  \item The last cell of s5 (OpenFHE CKKS 16384/1 relinearization) ran after an operating-system update, on kernel 7.0.0-42 instead of 7.0.0-38; its SEAL partner ran on 7.0.0-38. Binaries, images and configuration were unchanged (identical manifest). Its ratio (\R{CKKS:16384:1:relinearize}) is reported here but is \textbf{provisional}: the pair (SEAL and OpenFHE CKKS 16384/1 relinearization) will be measured again in one session, under one kernel, with the same protocol; the current cells stay in the raw data.
  \item The re-measured s4 OpenFHE addition ran at 4.27~GHz against 4.36--4.46~GHz for the rest of s4. It is an end-to-end API cost of a separate-analysis configuration, so no ratio depends on it.
  \item The per-second frequency record of the first parts of s2 and s3 was overwritten by a logging bug (fixed before s4); the measurements themselves are complete.
\end{itemize}

\begin{table}[h!]\centering\footnotesize\setlength{\tabcolsep}{4pt}
\input{\tab{conditions_s4s5.tex}}
\caption{Per-cell conditions in s4 and s5: mean effective core frequency and temperature during the cell, SEAL's ratio to the archived campaign, kernel.}
\label{tab:cond}
\end{table}

\clearpage
\section{Interpretation of the main differences}
\label{sec:interp}

Matched-profile comparisons are not exact-parameter comparisons: the primes, error distributions, CKKS scales (in some configurations) and key-switching methods still differ, as listed per cell. The observations below therefore hold for the stated configurations, boundaries, library versions and machine; none is a statement that one library is faster in general. Each item separates what was \emph{observed} from what would \emph{explain} it; no explanation below was tested.

\begin{enumerate}[leftmargin=*, itemsep=3pt]
  \item \textbf{Initialization and key generation.} The timed boundary is the same in both libraries: creation of the context from the parameters (in SEAL, including the generation of the primes), secret and public keys and, when the depth is positive, the relinearization key; no rotation keys. It is an initialization and key-generation cost, not the cost of an isolated key-generation primitive. \emph{Observed:} OpenFHE's cost was lower at CKKS 2048/5 (ratio \R{CKKS:2048:5:keygen}), 8192/1 (\R{CKKS:8192:1:keygen}) and 16384/1 (\R{CKKS:16384:1:keygen}); the per-library values of the separate-analysis configurations go in the same direction (Table~\ref{tab:separate}), without a ratio. \emph{Not established:} which part of the boundary (context creation, prime generation, key generation) causes the difference; it was not measured separately.
  \item \textbf{Decryption up to SEAL's boundary.} \emph{Observed:} ratio \ratioMinckksdecrypt--\ratioMaxckksdecrypt\ in all \ratioNckksdecrypt\ matched CKKS configurations (2048 to 16384, three sessions), \R{BFV:8192:5:decrypt} at BFV 8192/5 and \R{BGV:8192:5:decrypt} at BGV 8192/5. The CKKS ratio is the same where the initial scales are equal (8192/1: \R{CKKS:8192:1:decrypt}; 16384/1: \R{CKKS:16384:1:decrypt}) and where they differ, which is consistent with the scale not affecting this cost (established by test only at 8192/1). \emph{Not established:} the reason for the CKKS difference.
  \item \textbf{Multiplication without relinearization.} \emph{Observed:} ratio \R{CKKS:4096:1:multiply} at CKKS 4096/1, \R{CKKS:8192:1:multiply} at 8192/1 and \R{CKKS:16384:1:multiply} at 16384/1: close to parity, increasing with $N$ in these three configurations.
  \item \textbf{Key switching.} \emph{Observed:} at CKKS 8192/1, relinearization ratio \R{CKKS:8192:1:relinearize} and rotation ratio \R{CKKS:8192:1:rotate} (\RM{CKKS:8192:1:rotate} from medians), both reproduced in three independent sessions (Table~\ref{tab:repro}); at CKKS 16384/1, relinearization ratio \R{CKKS:16384:1:relinearize}, provisional until the pair is measured again under one kernel. This is the largest difference among the equivalent boundaries. \emph{Hypothesis, not tested:} the two libraries use different key-switching methods (SEAL: one special prime; OpenFHE: hybrid key switching, one prime per digit), so the difference may come from the method, the implementation, or both; the measurements do not separate them.
  \item \textbf{End-to-end API costs} (Table~\ref{tab:api}) are each library's cost for its own call, which includes different work, and are not compared. \emph{Observed within OpenFHE:} \texttt{EvalAdd}, which copies its operand, took \copyMin--\copyMax\ times as long as \texttt{EvalAddInPlace}; \texttt{Decrypt} with decoding took \V{OpenFHE:CKKS:16384:1:decrypt_full}~ms at CKKS 16384/1, against \V{OpenFHE:CKKS:16384:1:decrypt}~ms up to SEAL's boundary.
  \item \textbf{Separate-analysis cells} (all BFV and BGV cells except decryption at 8192/5; all of CKKS 8192/3; key generation and relinearization at several other CKKS configurations) are reported per library only. Their values show why no ratio is formed: at 8192/3, for example, OpenFHE's BGV relinearization took less time than SEAL's (\V{OpenFHE:BGV:8192:3:relinearize} against \V{SEAL:BGV:8192:3:relinearize}~ms), the reverse of the matched-profile CKKS cells, because the two libraries' moduli (number and size of primes) are different there.
\end{enumerate}

\section{Limitations and next steps}
\label{sec:limits}

\begin{itemize}[leftmargin=*, itemsep=2pt]
  \item \textbf{Latency only.} Energy and memory were not measured in this campaign; the plan is the SEAL campaign's dedicated protocol (timed-window RAPL with idle correction; memory at inner loop 1), for both libraries.
  \item \textbf{Frequency not fixed.} The CPU frequency followed the firmware's power management (boost enabled), which varied between days (Section~4). Comparisons are within a session, where conditions were close, and per-cell frequency is reported. For the next campaigns I propose to fix the frequency (boost disabled), so that results do not depend on the power state; this will be tested on this laptop first.
  \item \textbf{Scale equivalence} (CKKS) was established only for 8192/1 at $2^{30}$ and $2^{40}$. The initial scales are equal at 8192/1 and 16384/1; at 2048, 4096 and 8192/5 they differ by 2 to 14 bits, which is listed in the remaining differences of those cells (their ratios are decryption, key generation at 2048/5 and multiplication at 4096/1).
  \item \textbf{HElib and Lattigo} are not started, as agreed (SEAL first, then OpenFHE).
\end{itemize}

\section{Reproduction}

{\footnotesize
\begin{tabular}{@{}L{4.0cm}L{12.5cm}@{}}
\toprule
Step & Command (from the repository root) \\
\midrule
Build & \texttt{bash experiments/seal/docker\_build\_and\_run.sh}; \texttt{bash experiments/openfhe/docker\_build.sh} \\
One session (s1--s5) & \texttt{sudo bash experiments/tools/power\_guard.sh s1} (or directly \texttt{sudo bash experiments/openfhe/run\_standard\_campaign.sh s1}) \\
Statistics & \texttt{cd experiments/scripts; python3 aggregate.py --scenario=openfhe\_standard} \\
Tables of this note & \texttt{cd experiments/openfhe; python3 make\_standard\_tables.py} \\
Raw data, incidents & \texttt{experiments/results/raw/standard\_campaign/s1..s5/}, \texttt{INCIDENTS.md} \\
\bottomrule
\end{tabular}}

\end{document}
